% BHU PYQ -- Manually transcribed & error-corrected
% Paper: BCF-212 : Corporate Financial Management
% Exam: B.Com. (Hons.) F.M.M. Semester III Examination, 2023-24

\documentclass[11pt,a4paper]{article}
\usepackage[a4paper,top=2.5cm,bottom=2.5cm,left=2.8cm,right=2.8cm,headheight=14pt]{geometry}
\usepackage{amsmath,amssymb}
\usepackage[T1]{fontenc}
\usepackage[utf8]{inputenc}
\usepackage{lmodern,microtype,fancyhdr,enumitem,hyperref,lastpage,parskip,booktabs}

\hypersetup{
    colorlinks=true,
    urlcolor=black,
    linkcolor=black,
    pdftitle={BCF-212: Corporate Financial Management -- BHU B.Com. (Hons.) FMM Sem III 2023-24}
}

\pagestyle{fancy}
\fancyhf{}
\renewcommand{\headrulewidth}{0.4pt}
\renewcommand{\footrulewidth}{0.4pt}
\fancyhead[L]{\small   \textit{Banaras Hindu University}}
\fancyhead[R]{\small   \textit{BCF-212: Corporate Financial Management}}
\fancyfoot[C]{\small Page  \thepage\ of \pageref{LastPage}}


\newlist{parts}{enumerate}{2}
\setlist[parts,1]{label=(\alph*),leftmargin=*,itemsep=5pt,topsep=3pt}
\setlist[parts,2]{label=(\roman*),leftmargin=*,itemsep=3pt,topsep=2pt}

\newcommand{\pts}[1]{\hfill   \textit{[#1]}}


\begin{document}


\begin{center}
  {\large\bfseries BANARAS HINDU UNIVERSITY}\\[2pt]
  {\normalsize B.Com. (Hons.) F.M.M. --- Semester III Examination, 2023-24}\\[6pt]
  \rule{0.8\linewidth}{0.5pt}\\[6pt]
  {\large\bfseries Paper No. BCF-212: Corporate Financial Management}\\[4pt]
  {\normalsize Time: 3 Hours\hfill Maximum Marks: 70}
\end{center}

\rule{\linewidth}{0.4pt}

\smallskip



\noindent   \textit{Instructions: Answer all questions. All questions carry equal marks (14 marks each).}

\bigskip




\noindent   \textbf{Question 1.}
``The goal of profit maximization does not provide an operationally useful criterion for measuring the success of business operations.'' Explain the statement.\pts{14}

\centerline{   \textbf{OR}}

``Financial management, as an integral part of the overall management, is not a totally independent area.'' Justify.\pts{14}

\medskip\rule{\linewidth}{0.2pt}




\noindent   \textbf{Question 2.}
Discuss the factors that influence the requirements of fixed capital of a company.\pts{14}

\centerline{   \textbf{OR}}

Explain the characteristics of a sound financial plan. Also describe the types of financial plans.\pts{14}

\medskip\rule{\linewidth}{0.2pt}




\noindent   \textbf{Question 3.}
Critically examine the advantages and disadvantages of raising funds by issuing shares of different types.\pts{14}

\centerline{   \textbf{OR}}

Examine the utility of public deposits as a source of corporate finance in India.\pts{14}

\medskip\rule{\linewidth}{0.2pt}




\noindent   \textbf{Question 4.}
What do you mean by over-capitalisation? Discuss the causes and remedies of over-capitalisation.\pts{14}

\centerline{   \textbf{OR}}

Distinguish between capitalization and capital structure.\pts{14}

\medskip\rule{\linewidth}{0.2pt}




\noindent   \textbf{Question 5.}
Explain the `Net Operating Income' approach of Capital Structure.\pts{14}

\centerline{   \textbf{OR}}

Discuss briefly the factors governing the design of capital structure.\pts{14}

\vfill
\rule{\linewidth}{0.4pt}

\begin{center}\small
\href{https://bhu-exam.vercel.app/}{Click here to visit our website}\\[4pt]
\href{https://me-aryan.vercel.app/}{Click here to visit our contributor}\\[4pt]
\href{https://quashan-me.vercel.app/}{Click here to visit our contributor}
\end{center}

\end{document}
